\section*{Chapter \ref{ch:iv1:number-riddles} --- Number Riddles}
\markboth{\protect\omTocMark{}{Answers: Number Riddles}}{\protect\omTocMark{}{Answers: Number Riddles}}

\begin{solution}{iq:iv1:number-riddles:1}
The angle is 78 degrees.

Work from 12 o'clock. The minute hand moves 6 degrees a minute, so at 36 minutes it is at 216 degrees. The hour hand moves half a degree a minute, so at 4:36 it is at $4 \times 30 + 36 \times 0.5 = 138$ degrees. The difference is 78 degrees.

\iqwrong{Using $120$ degrees for the hour hand, as if it stayed on the 4 for the whole hour. That gives 96.}
\iqfollow{When do the hands first coincide after 4:00?}{The minute hand gains $5.5$ degrees a minute, and it starts 120 degrees behind, so after $120/5.5 = 21.82$ minutes, which is about 4:21:49.}
\iqfollow{How often do they coincide?}{Eleven times in 12 hours, so every $12/11$ of an hour, which is every 65 minutes and about 27 seconds.}
\end{solution}

\begin{solution}{iq:iv1:number-riddles:2}
$250!$ ends in 62 zeros.

A trailing zero needs a factor of 10, which is a 2 and a 5. There are many more 2s than 5s among the factors, so the count of 5s decides it. Using Legendre's formula\index{Legendre's formula}, multiples of 5 give $250/5 = 50$, multiples of 25 give a further 10, and multiples of 125 give a further 2. The total is 62.

\iqwrong{Stopping at $250/5 = 50$. A multiple of 25 holds two 5s and a multiple of 125 holds three.}
\iqfollow{Which factorial ends in exactly 5 zeros?}{None. $24!$ has 4 zeros, and $25!$ has 6, because the factor 25 brings two 5s at once.}
\iqfollow{And the last non-zero digit?}{That is harder, since you must track the unmatched 2s. For $10!$, the number is 3628800, so the last non-zero digit is 8.}
\end{solution}

\begin{solution}{iq:iv1:number-riddles:3}
The last digit is 9.

Only last digits matter for the last digit of a product. Powers of 7 end in 7, 9, 3, 1 and then repeat, with period 4. Since $2026 = 4 \times 506 + 2$, the power has the same last digit as $7^2 = 49$.

\iqwrong{Using 2026 as if the last digit cycled with period 10. Look at the first few powers, and read the period off directly.}
\iqfollow{What is the last digit of $7^{7^7}$?}{You need the exponent modulo 4. Since $7 \equiv -1$ modulo 4 and $7^7$ is an odd power, the exponent is $\equiv 3$ modulo 4. So the last digit is that of $7^3 = 343$, which is 3.}
\iqfollow{Which digits have a period shorter than 4?}{The digits 0, 1, 5 and 6 stay fixed, and 4 and 9 have period 2. The digits 2, 3, 7 and 8 have period 4.}
\end{solution}

\begin{solution}{iq:iv1:number-riddles:4}
31 doors are open, and door 72 is changed 12 times.

Door $n$ is changed once for each divisor of $n$. Divisors come in pairs, $d$ and $n/d$, so the count is even unless the two members of a pair coincide, which happens only for a perfect square. So the open doors are the perfect squares up to 1000: $1^2$ up to $31^2 = 961$, since $32^2 = 1024$.

For door 72, write $72 = 2^3 \times 3^2$. A divisor chooses a power of 2 from 4 options and a power of 3 from 3, so there are $4 \times 3 = 12$. Being even, door 72 ends closed.

\iqwrong{Guessing that half the doors end open, as if each one were a coin toss. The number of changes is entirely decided by the door's divisors.}
\iqfollow{Which door below 1000 is changed most often?}{Door 840, which has 32 divisors, more than any other number below 1000. It ends closed.}
\end{solution}

\begin{solution}{iq:iv1:number-riddles:5}
Writing a number $n$ twice gives $1001 \times n$, and $1001 = 7 \times 11 \times 13$.

Writing $n$ three places to the left and adding $n$ is the same as $1000 n + n$. For example, $123123 = 123 \times 1001$. Dividing by 7, then 11, then 13 gives back $n$.

\iqfollow{How does this give a test for divisibility by 7?}{Since $1000 \equiv -1$ modulo 1001, and so modulo 7, split a number into blocks of three digits from the right, and take the alternating sum. For $123456$ the sum is $456 - 123 = 333$, which leaves remainder 4 on division by 7, and so does $123456$. The same test works for 11 and 13.}
\iqfollow{What does writing a two-digit number twice do?}{It multiplies by 101, which is prime. So $4747 = 47 \times 101$.}
\iqwrong{Trying examples until you see a pattern. Examples will convince you, but the question asks why, and only the factor of 1001 explains it.}
\end{solution}

\begin{solution}{iq:iv1:number-riddles:6}
The number is 2519.

In each case the remainder is one less than the divisor. So the number plus 1 is divisible by every $k$ from 2 to 10, and the smallest such number plus 1 is the least common multiple\index{least common multiple}: $\mathrm{lcm}(2, \dots, 10) = 2^3 \times 3^2 \times 5 \times 7 = 2520$. So the number is 2519.

\iqwrong{Answering $10! - 1 = 3\,628\,799$. That does satisfy every condition, but it is not the smallest: the factorial counts repeated factors that the least common multiple does not need.}
\iqfollow{And up to 20?}{The least common multiple adds the primes 11, 13, 17 and 19 and one more power of 2: $2^4 \times 3^2 \times 5 \times 7 \times 11 \times 13 \times 17 \times 19 = 232\,792\,560$. So the number is $232\,792\,559$.}
\end{solution}

\begin{solution}{iq:iv1:number-riddles:7}
$20!$ ends in 8 zeros in base 12.

A trailing zero in base 12 is a factor of 12, which is $2^2 \times 3$. Count both primes with Legendre's formula\index{Legendre's formula}. For 2: $10 + 5 + 2 + 1 = 18$. For 3: $6 + 2 = 8$. Each factor of 12 uses two 2s and one 3, so 2s allow $18/2 = 9$ factors of 12, and 3s allow 8. The smaller count wins, so 8 zeros.

\iqwrong{Trusting the base-10 habit of counting only the larger prime. It happens to give 8 here, but you must compare the count of each prime against its exponent in 12, because the limiting prime is not always the larger one. For example, $9!$ has seven 2s and four 3s, so it ends in 3 zeros in base 12, limited by the 2s, and not in 4.}
\iqfollow{What about base 6 and base 10?}{Base 6 needs a 2 and a 3 per zero, so $\min(18, 8) = 8$. Base 10 needs a 2 and a 5, so $\min(18, 4) = 4$.}
\end{solution}

\begin{solution}{iq:iv1:number-riddles:8}
$a = 3$ and $b = 2$, so the number is 43272.

Split 72 into its coprime factors, 8 and 9, and use a separate test for each. Divisibility by 8 depends only on the last three digits, $27b$. Among 270 to 279, only 272 is a multiple of 8, so $b = 2$. Divisibility by 9 depends on the digit sum, $4 + a + 2 + 7 + 2 = 15 + a$, which must be a multiple of 9, so $a = 3$. Check: $43272 = 72 \times 601$.

\iqwrong{Testing multiples of 72, of which there are too many. Use only coprime factors: 72 is $8 \times 9$, not $6 \times 12$, because a number can be a multiple of both 6 and 12 without being a multiple of 72.}
\iqfollow{Is the answer always unique?}{No. For $7a36b$, $b$ may be 0 or 8, and these give 72360 and 73368. Check every candidate before you stop.}
\end{solution}

\enlargethispage{-3\baselineskip}%
\begin{solution}{iq:iv1:number-riddles:9}
The remainder is 8.

Use the small powers. Modulo 17, $3^2 = 9$, $3^4 = 81 \equiv 13 \equiv -4$, and $3^8 \equiv 16 \equiv -1$. So $3^{16} \equiv 1$, which is what Fermat's little theorem\index{Fermat's little theorem} promises for any number not divisible by 17. Now $2026 = 16 \times 126 + 10$, so
\[
3^{2026} \equiv 3^{10} = 3^8 \times 3^2 \equiv -1 \times 9 = -9 \equiv 8.
\]

\iqwrong{Trying to reduce the exponent modulo 17. The period of powers divides $17 - 1 = 16$, so you reduce the exponent modulo 16.}
\iqfollow{Does every base have period 16 here?}{No. It divides 16. Modulo 17, 3 has period 16, but $2^4 = 16 \equiv -1$, so 2 has period 8.}
\iqwhy{Fast exponentiation and modular arithmetic of this kind are behind hash functions and random-number generators, so a developer may well be asked to code this after working it by hand.}
\end{solution}

\begin{solution}{iq:iv1:number-riddles:10}
The smallest is 301, and the next is 721.

The first five conditions say that $n - 1$ is a multiple of $\mathrm{lcm}(2, 3, 4, 5, 6) = 60$, so $n = 60k + 1$. Now apply the last condition, modulo 7. Since $60 \equiv 4$, you need $4k + 1 \equiv 0$, which holds at $k = 5$, because $4 \times 5 + 1 = 21$. So $n = 301$, which is $7 \times 43$.

The solutions repeat every $\mathrm{lcm}(60, 7) = 420$. That is the Chinese remainder theorem\index{Chinese remainder theorem}: with coprime moduli, there is exactly one solution in each full period. The next is therefore 721.

\iqwrong{Using $2 \times 3 \times 4 \times 5 \times 6 = 720$ as the period, which gives 721 as the smallest. The common multiple is 60, not 720, and 721 is only the second solution.}
\iqfollow{Solve remainders 2, 3, 2 on division by 3, 5, 7.}{Numbers with remainder 2 for both 3 and 7 are $21k + 2$. Modulo 5 this is $k + 2 \equiv 3$, so $k = 1$, and the answer is 23.}
\end{solution}

\begin{solution}{iq:iv1:number-riddles:11}
The largest amount you cannot pay is 59, and there are 30 amounts you cannot pay.

For two coins $a$ and $b$ with no common factor, the largest unreachable amount is $ab - a - b$, which is $77 - 18 = 59$. The count of unreachable amounts is $(a - 1)(b - 1)/2 = 6 \times 10 / 2 = 30$.

Use remainders modulo 7\index{residue argument}. Using 7s alone moves you up within a remainder class. So in each class you want to know the smallest amount reachable with 11s: these are 0, 11, 22, 33, 44, 55 and 66, one for each remainder. Once a class has been reached, every later amount in it is reachable. The last amount to be missed is just below the largest of these: $66 - 7 = 59$. Above that, every class has started.

\iqfollow{And with coins worth 6 and 9?}{Their common factor is 3, so no amount that is not a multiple of 3 is ever reachable, and there is no largest unreachable amount.}
\iqwrong{Multiplying and answering $77 - 1$, or thinking that 77 is the first point past which all amounts work. The coins start to cover every class well before 77.}
\end{solution}

\begin{solution}{iq:iv1:number-riddles:12}
You write the digit 1 exactly 301 times, and the digits from 1 to 999 sum to 13\,500.

Treat the numbers 0 to 999 as three-digit strings with leading zeros\index{symmetry arguments}. By symmetry, each digit appears equally often in each place. Each place holds each digit 100 times, so the digit 1 appears $3 \times 100 = 300$ times. The number 1000 adds one more, for 301.

For the sum, each place contributes $100 \times (0 + 1 + \dots + 9) = 4500$, and three places give $13\,500$. The leading zeros and the number 0 add nothing to it, so including them is harmless.

\iqwrong{Counting how many numbers contain a 1. That is a different question: among 0 to 999, the numbers without a 1 number $9^3 = 729$, so 271 contain a 1. With 1000, this makes 272, while a number such as 111 still contributes three digits to the count of 301.}
\iqfollow{And the digit 0, from 1 to 1000?}{Now the leading zeros do matter. The answer is 192: 9 in the numbers 10 to 90, 180 in the three-digit numbers, and 3 in 1000.}
\end{solution}

\begin{solution}{iq:iv1:number-riddles:13}
With either pan allowed, four weights are enough: 1, 3, 9 and 27. With one pan, you need six: 1, 2, 4, 8, 16 and 32.

With both pans, each weight has three states: on the object's pan, on the other pan, or off. So $n$ weights give $3^n$ signed sums. These are symmetric about 0, so at most $(3^n - 1)/2$ positive values are reachable. For $n = 3$ that is 13, which is too few, and for $n = 4$ it is 40, exactly enough. The powers of 3 achieve it, as every number has one way of being written in balanced ternary\index{balanced ternary}, with digits $-1$, 0 and 1. For example, $23 = 27 - 3 - 1$, and $40 = 27 + 9 + 3 + 1$.

With one pan, each weight is on or off, so you are writing the number in binary\index{number bases}. Five weights reach only 31, and six reach 63.

\iqwrong{Answering 1, 2, 4, 8, 16, 32 for both versions. Binary is optimal for one pan, but it wastes the third state.}
\iqfollow{Which weights are used for 5 grams?}{$5 = 9 - 3 - 1$: put the 9 on the opposite pan from the object, and the 3 and 1 with the object.}
\end{solution}

\begin{solution}{iq:iv1:number-riddles:14}
The last two digits are 01.

Work modulo 100 by repeated squaring. $3^4 = 81$, so $3^5 = 243 \equiv 43$. Then $3^{10} \equiv 43^2 = 1849 \equiv 49$, and $3^{20} \equiv 49^2 = 2401 \equiv 1$. So the powers of 3 repeat with period 20 modulo 100. Since $1000 = 20 \times 50$, $3^{1000} \equiv 1$.

The general reason is Euler's theorem\index{Euler's theorem}: a number coprime to 100 has period dividing $\varphi(100) = 40$. Here the true period is 20, which is shorter.

\iqwrong{Thinking that the theorem applies to every base. It needs the base to share no factor with 100.}
\iqfollow{So what are the last two digits of $2^{1000}$?}{76. The base 2 shares a factor with 100, so split 100 into 4 and 25. Modulo 4 the power is 0. Modulo 25 the period divides 20, so the power is 1. The only number below 100 that is a multiple of 4 and one more than a multiple of 25 is 76.}
\end{solution}

\begin{solution}{iq:iv1:number-riddles:15}
The highest power of 2 is $2^3$, the highest power of 3 is $3^4$, and 7 does not divide it.

Write $\binom{100}{50} = 100! / (50!)^2$, and count each prime with Legendre's formula\index{Legendre's formula}. For 2: $100!$ holds $50 + 25 + 12 + 6 + 3 + 1 = 97$, and $50!$ holds $25 + 12 + 6 + 3 + 1 = 47$, so the answer is $97 - 94 = 3$. For 3: $48$ in $100!$ and $22$ in $50!$, which leaves $48 - 44 = 4$. For 7: $14 + 2 = 16$ in $100!$ and $7 + 1 = 8$ in $50!$, which leaves 0.

Kummer's theorem\index{Kummer's theorem} says this another way: the power of a prime $p$ in $\binom{a + b}{a}$ equals the number of carries when adding $a$ and $b$ in base $p$. In base 7, 50 is $101$, and adding it to itself gives $202$ with no carries, so 7 does not divide it. In base 2, $50 = 110010$ and there are three carries.

\iqfollow{Why must the number be even?}{$\binom{2n}{n}$ is always even, as $\binom{2n}{n} = 2 \binom{2n - 1}{n - 1}$.}
\iqwrong{Assuming that every prime up to 100 must divide the number, since $100!$ contains them all. The denominator can cancel every copy, as it does here for 7.}
\end{solution}

\enlargethispage{-3\baselineskip}%
\begin{solution}{iq:iv1:number-riddles:16}
Factor: $p^2 - 1 = (p - 1)(p + 1)$, and show that 8 and 3 divide this separately.

Since $p$ is odd, $p - 1$ and $p + 1$ are consecutive even numbers. One of them is a multiple of 4, and the other a multiple of 2, so the product is a multiple of 8. Next, among the three consecutive numbers $p - 1$, $p$ and $p + 1$ one is a multiple of 3. It is not $p$, as $p$ is a prime greater than 3, so 3 divides the product. As 8 and 3 are coprime, 24 divides it. For $p = 5$, $p^2 - 1 = 24$, and for $p = 7$ it is 48, so 24 cannot be replaced by anything larger.

\iqfollow{Is primality really needed?}{No. It uses only that $p$ is coprime to 6, so $25^2 - 1 = 624 = 24 \times 26$ works as well.}
\iqwrong{Testing 5, 7 and 11 and declaring it proved. Examples show that a claim could be true, and a proof must cover every case.}
\iqfollow{Is there a similar statement for fourth powers?}{Yes. For a prime above 5, $p^4 - 1$ is a multiple of 240, by the same reasoning with 16, 3 and 5.}
\end{solution}

\begin{solution}{iq:iv1:number-riddles:17}
The number is 2178, since $4 \times 2178 = 8712$.

Write it as $abcd$, with $4 \times abcd = dcba$. The product has four digits, so $a$ is 1 or 2. The product's last digit is $a$ and it is also the last digit of $4d$, which is even, so $a = 2$. Then $d \geq 8$, and $4d$ must end in 2, which rules out $d = 9$ and leaves $d = 8$.

Now expand with $a = 2$ and $d = 8$:
\[
4(2008 + 100b + 10c) = 8002 + 100c + 10b.
\]
This simplifies to $13b + 1 = 2c$. Since $c \leq 9$, $b$ is at most 1. And $b = 0$ would make $2c$ odd, so $b = 1$ and $c = 7$. So the number is 2178.

\iqfollow{What is the matching number for 9?}{1089, since $9 \times 1089 = 9801$. Putting a 9 in the middle continues the pattern: $4 \times 21978 = 87912$, and $9 \times 10989 = 98901$.}
\iqwrong{Searching through thousands of numbers. The first and last digits pin down $a$ and $d$ in seconds, and the rest is one line of algebra.}
\end{solution}

\begin{solution}{iq:iv1:number-riddles:18}
The sum is 360.

Pair the numbers by symmetry\index{symmetry arguments}. For $0.abcabc\dots$ take its partner with digits $9 - a$, $9 - b$ and $9 - c$. The partner's digits are also different from one another. The two numbers add to $0.999\dots = 1$. No number is its own partner, as $a = 9 - a$ has no whole-number solution.

The count of such numbers is $10 \times 9 \times 8 = 720$, which makes 360 pairs, each adding to 1. So the total is 360.

You can check it another way. The value of $0.abcabc\dots$ is $abc/999$. Each digit appears in each place 72 times, and on average the three-digit string is $499.5$, so the total is $720 \times 499.5 / 999 = 360$.

\iqwrong{Trying to sum in the order of the digits. It can be done, but the pairing removes all the arithmetic. If a sum has a symmetry, look for it first.}
\iqfollow{What about all $0.abab\dots$ with $a \neq b$?}{There are 90, and the same pairing with $9 - a$ gives 45 pairs. The sum is 45.}
\end{solution}

\begin{solution}{iq:iv1:number-riddles:19}
There are 132 such moments, in 66 pairs.

Measure time $t$ in units of 12 hours, so $0 \le t < 1$. The hour hand is at $t$ of a full turn, and the minute hand at $12t$, modulo 1. A reading is ambiguous if swapping the hands also gives a valid time $s$. That means $s = 12t$ and $t = 12s$, modulo 1. Substituting, $144 t \equiv t$, so $143 t$ is a whole number. There are 143 such values of $t$: $k/143$ for $k = 0, \dots, 142$.

Among these, some have the hands together, so that the swap gives the same time. These are the coincidences, with $11t$ a whole number, and there are 11 of them. They are not ambiguous. That leaves $143 - 11 = 132$. Each pairs with a different time that shows the same picture, so they form 66 pairs.

\iqwrong{Answering 143. The 11 moments when the hands overlap do show a unique time, since swapping two identical hands in the same place changes nothing.}
\iqfollow{Is the moment $t = 1/143$ a real example?}{Yes. It is 5 minutes and about 2 seconds past 12. The hour hand is at $1/143$ of a turn and the minute hand at $12/143$. Its swap is $t = 12/143$, which is 1:00:25 or so, where the hour hand is at $12/143$ and the minute hand at $144/143 \equiv 1/143$.}
\end{solution}

\begin{solution}{iq:iv1:number-riddles:20}
The largest is 43, and every number above it can be bought.

Work modulo 6\index{residue argument}. Boxes of 6 let you move up within a remainder class, so you only need the smallest amount in each class that the 9s and 20s can reach:
\[
\begin{array}{c|cccccc}
\text{remainder} & 0 & 1 & 2 & 3 & 4 & 5 \\ \hline
\text{smallest amount} & 0 & 49 & 20 & 9 & 40 & 29
\end{array}
\]
Here $49 = 9 + 40$, $40 = 20 + 20$ and $29 = 9 + 20$. Each is the smallest because reaching a remainder of 1, 2, 4 or 5 needs at least one box of 20, and each extra box of 9 only moves you by 3 modulo 6. So the last missing number is in the class of 49, just below it: $49 - 6 = 43$. Every number from 44 up lies at or above the smallest amount of its class, so 6s finish it.

\iqwrong{Using the two-coin formula on 9 and 20 and answering 151. That is right if no boxes of 6 exist. There is no such simple formula for three coin values\index{Frobenius coin problem}, so you must do the table.}
\iqfollow{What does 44 look like, as a check?}{$44 = 20 + 6 \times 4$, and $45 = 9 \times 5$, $46 = 20 + 20 + 6$, $47 = 20 + 9 \times 3$, $48 = 6 \times 8$, $49 = 20 + 20 + 9$. These six in a row, plus 6s, cover all that follow.}
\end{solution}

\begin{solution}{iq:iv1:number-riddles:21}
The powers of 2 are exactly the numbers that cannot be written this way. And 2026 can be written in exactly one way: $505 + 506 + 507 + 508$.

Suppose that $k \geq 2$ consecutive numbers starting at $m \geq 1$ add to $n$. Their sum is $k(2m + k - 1)/2$, so
\[
2n = k \, (2m + k - 1).
\]
The two factors have opposite parity, as their difference is $2m - 1$, which is odd. The second is larger than $k$. If $n$ is a power of 2 then so is $2n$, which has no odd factor above 1. So the odd one of the two factors must be 1. But $k \geq 2$, and the other factor is larger still, so this is impossible.

Otherwise, $n$ has an odd factor $d > 1$. Write $2n$ as a product of the odd part $d$ and the even part $2n/d$. Call the smaller one $k$ and the larger $2m + k - 1$; they have opposite parity, so $m$ is a whole number, and being the larger makes $m \ge 1$. Each odd divisor above 1 gives exactly one way. As $2026 = 2 \times 1013$ with 1013 prime, there is one way, with $k = 4$ and $2m + 3 = 1013$.

\iqfollow{How many ways for 15?}{Its odd divisors above 1 are 3, 5 and 15, so three ways: $7 + 8$, $4 + 5 + 6$ and $1 + 2 + 3 + 4 + 5$.}
\end{solution}

\enlargethispage{-3\baselineskip}%
\begin{solution}{iq:iv1:number-riddles:22}
Person 73 is left.

First take the case where the number of people is a power of 2. Each full pass round the circle removes half the people, and ends just before person 1, so the count starts from person 1 again. So person 1 is left when the count is a power of 2.

Now $100 = 64 + 36$. After 36 people have gone, 64 remain, and the count has reached the person just after the 36th to leave. The 36th to leave is person 72, so the next to stay is person 73, and from there it is the power-of-2 case with person 73 in the role of person 1.

In general, for $n = 2^a + b$ with $0 \le b < 2^a$, the survivor is $2b + 1$\index{Josephus problem}. In binary this moves the leading digit of $n$ to the end: $100$ is $1100100$, and that gives $1001001$, which is 73.

\iqwrong{Answering 1 by the power-of-2 argument, which holds only when the count is a power of 2. Here person 1 survives the first two passes. With 25 left the pass ends on a person who stays, and so person 1 leaves first in the next pass.}
\iqfollow{What if every third person leaves instead?}{The neat form is lost. Use the recurrence $J(n) = (J(n - 1) + 3) \bmod n$, with people counted from 0. That reaches the number 91 for 100 people, in the numbering of this question.}
\end{solution}

\begin{solution}{iq:iv1:number-riddles:23}
The number is $105\,263\,157\,894\,736\,842$. It has 18 digits, and doubling it gives $210\,526\,315\,789\,473\,684$.

Let the number be $N = 10a + d$, where $d$ is the last digit and $a$ has $n - 1$ digits. Moving $d$ to the front gives $d \cdot 10^{n-1} + a$, and this must equal $2N = 20a + 2d$. So
\[
19a = d \,(10^{n-1} - 2).
\]
Since $d < 19$ and 19 is prime, 19 must divide $10^{n-1} - 2$. So we need the smallest $m = n - 1$ with $10^m \equiv 2$ modulo 19. Working through the powers of 10 gives $m = 17$.

Then $a = d(10^{17} - 2)/19$. With $d = 1$ this has only 16 digits, which would force a leading zero in $N$. The smallest valid value is $d = 2$, which gives $a = 10\,526\,315\,789\,473\,684$ with all 17 digits.

You can also build it by hand from the right: start with last digit 2, and double repeatedly, writing each digit of the product to the left of the last. You stop after writing a 1 with no carry out, because its double is 2, the digit you started from. That is the quickest way to the digits by hand.

\iqfollow{Why are these digits familiar?}{They are the repeating block of $2/19$ written as a decimal.}
\iqfollow{And for tripling?}{The same steps with $29a = d(10^{n-1} - 3)$ give a 28-digit answer, $1\,034\,482\,758\,620\,689\,655\,172\,413\,793$.}
\end{solution}

\begin{solution}{iq:iv1:number-riddles:24}
$C = 7$.

The trick is not to compute the number. Two facts settle it: the sum of digits keeps the remainder modulo 9, and the digit sums shrink very fast.

First, the remainder. $4444 \equiv 16 \equiv 7$ modulo 9. Powers of 7 modulo 9 go 7, 4, 1 and repeat, with period 3. Since $4444 \equiv 1$ modulo 3, $4444^{4444} \equiv 7$ modulo 9, and so $C \equiv 7$ as well\index{casting out nines}.

Next, the size. The number is below $10\,000^{4444} = 10^{17\,776}$, so it has at most $17\,776$ digits, and $A \le 9 \times 17\,776 = 159\,984$. A number up to that size has six digits at most, and the first of six is 1, so $B \le 1 + 5 \times 9 = 46$. Up to 46, the largest digit sum belongs to 39, so $C \le 12$. The only number up to 12 with remainder 7 on division by 9 is 7.

\iqwrong{Assuming that the digit sum of a large power is about $4.5$ a digit and estimating from there. That gives $A$ and $B$ only roughly, and the remainder is what fixes the answer.}
\iqfollow{What are the actual values?}{Direct computation gives $A = 72\,601$, $B = 16$ and $C = 7$. Note that $A$ is well below its bound, near $4.5$ a digit.}
\end{solution}

\begin{solution}{iq:iv1:number-riddles:25}
The smallest possible number is 3121.

Each sailor takes a pile $x$ and leaves $\tfrac45(x - 1)$. The neat step is to add four ghost coconuts: since $\tfrac45(x - 1) + 4 = \tfrac45(x + 4)$, the quantity $x + 4$ simply gets multiplied by $\tfrac45$ at each step\index{invariant shift}. So after five sailors,
\[
\text{final} + 4 = \left(\tfrac45\right)^5 (N + 4).
\]
For this to be whole, $5^5$ must divide $N + 4$, so $N + 4 = 3125\,k$ and the final pile is $1024\,k - 4$. The final pile must also divide by 5, and $1024\,k - 4 \equiv 4k - 4$ modulo 5, so $k \equiv 1$. The smallest is $k = 1$, which gives $N = 3121$.

Check: the piles go $3121$, $2496$, $1996$, $1596$, $1276$, $1020$, and $1020 = 5 \times 204$.

The value $-4$ is the fixed point of the step, since $\tfrac45(-4 - 1) = -4$. If negative coconuts were allowed, a pile of $-4$ would stay $-4$ forever, and this is what suggests the ghost coconuts.

\iqwrong{Working backwards from a guess at the final pile and trying each in turn. It will reach 3121, but you must try many values, and you cannot tell whether you have the smallest.}
\iqfollow{And if one more coconut goes to the monkey in the morning?}{The final pile must be one more than a multiple of 5: $1024\,k - 4 \equiv 1$ needs $k \equiv 0$ modulo 5, so $k = 5$, which gives $N = 5^6 - 4 = 15\,621$.}
\end{solution}
